Bodový graf závislosti čistého příjmu domácností v~\ac{PPS} na pokrytí \ac{KS} je stěžejním grafem této práce: dokumentuje kladnou korelaci, podle níž státy s~vyšším pokrytím \ac{KS} vykazují vyšší reálné čisté příjmy --- a~to i~po kontrole na úroveň \ac{HDP} na obyvatele.
Kauzální interpretace tohoto vztahu je podpořena teorií vyjednávání: \ac{KS} s~plošnou platností komprimují mzdové rozdíly a~zvyšují spodní část mzdové distribuce, čímž přispívají k~vyšší průměrné a~mediánové úrovni příjmu v~celé ekonomice.
ČR se umísťuje nalevo od středu grafu s~relativně nízkým příjmem v~\ac{PPS} a~nízkým \ac{KS} pokrytím (cca \SI{35}{\percent}): hypotetické přiblížení k~pokrytí \SIrange{70}{80}{\percent} (AT úroveň) by dle korelační linie odpovídalo nárůstu \ac{PPS} příjmu o~\SIrange{15}{20}{\percent} --- efekt srovnatelný s~výstupy z~přirozených experimentů v~literatuře.
Tento obrázek může sloužit jako ústřední vizualizace teze: posílení \ac{KV} pokrytí není jen otázkou pracovního práva, ale klíčovým makroekonomickým nástrojem pro zvýšení životní úrovně bez nutnosti přímé státní redistribuce.

Bodový graf závislosti čistého příjmu domácností v~\ac{PPS} na pokrytí \ac{KS} je stěžejním grafem této práce: dokumentuje kladnou korelaci, podle níž státy s~vyšším pokrytím \ac{KS} vykazují vyšší reálné čisté příjmy --- a~to i~po kontrole na úroveň \ac{HDP} na obyvatele.
Kauzální interpretace tohoto vztahu je podpořena teorií vyjednávání: \ac{KS} s~plošnou platností komprimují mzdové rozdíly a~zvyšují spodní část mzdové distribuce, čímž přispívají k~vyšší průměrné a~mediánové úrovni příjmu v~celé ekonomice.
ČR se umísťuje nalevo od středu grafu s~relativně nízkým příjmem v~\ac{PPS} a~nízkým \ac{KS} pokrytím (cca \SI{35}{\percent}): hypotetické přiblížení k~pokrytí \SIrange{70}{80}{\percent} (AT úroveň) by dle korelační linie odpovídalo nárůstu \ac{PPS} příjmu o~\SIrange{15}{20}{\percent} --- efekt srovnatelný s~výstupy z~přirozených experimentů v~literatuře.
Tento obrázek může sloužit jako ústřední vizualizace teze: posílení \ac{KV} pokrytí není jen otázkou pracovního práva, ale klíčovým makroekonomickým nástrojem pro zvýšení životní úrovně bez nutnosti přímé státní redistribuce.

Bodový graf závislosti čistého příjmu domácností v~\ac{PPS} na pokrytí \ac{KS} je stěžejním grafem této práce: dokumentuje kladnou korelaci, podle níž státy s~vyšším pokrytím \ac{KS} vykazují vyšší reálné čisté příjmy --- a~to i~po kontrole na úroveň \ac{HDP} na obyvatele.
Kauzální interpretace tohoto vztahu je podpořena teorií vyjednávání: \ac{KS} s~plošnou platností komprimují mzdové rozdíly a~zvyšují spodní část mzdové distribuce, čímž přispívají k~vyšší průměrné a~mediánové úrovni příjmu v~celé ekonomice.
ČR se umísťuje nalevo od středu grafu s~relativně nízkým příjmem v~\ac{PPS} a~nízkým \ac{KS} pokrytím (cca \SI{35}{\percent}): hypotetické přiblížení k~pokrytí \SIrange{70}{80}{\percent} (AT úroveň) by dle korelační linie odpovídalo nárůstu \ac{PPS} příjmu o~\SIrange{15}{20}{\percent} --- efekt srovnatelný s~výstupy z~přirozených experimentů v~literatuře.
Tento obrázek může sloužit jako ústřední vizualizace teze: posílení \ac{KV} pokrytí není jen otázkou pracovního práva, ale klíčovým makroekonomickým nástrojem pro zvýšení životní úrovně bez nutnosti přímé státní redistribuce.

Bodový graf závislosti čistého příjmu domácností v~\ac{PPS} na pokrytí \ac{KS} je stěžejním grafem této práce: dokumentuje kladnou korelaci, podle níž státy s~vyšším pokrytím \ac{KS} vykazují vyšší reálné čisté příjmy --- a~to i~po kontrole na úroveň \ac{HDP} na obyvatele.
Kauzální interpretace tohoto vztahu je podpořena teorií vyjednávání: \ac{KS} s~plošnou platností komprimují mzdové rozdíly a~zvyšují spodní část mzdové distribuce, čímž přispívají k~vyšší průměrné a~mediánové úrovni příjmu v~celé ekonomice.
ČR se umísťuje nalevo od středu grafu s~relativně nízkým příjmem v~\ac{PPS} a~nízkým \ac{KS} pokrytím (cca \SI{35}{\percent}): hypotetické přiblížení k~pokrytí \SIrange{70}{80}{\percent} (AT úroveň) by dle korelační linie odpovídalo nárůstu \ac{PPS} příjmu o~\SIrange{15}{20}{\percent} --- efekt srovnatelný s~výstupy z~přirozených experimentů v~literatuře.
Tento obrázek může sloužit jako ústřední vizualizace teze: posílení \ac{KV} pokrytí není jen otázkou pracovního práva, ale klíčovým makroekonomickým nástrojem pro zvýšení životní úrovně bez nutnosti přímé státní redistribuce.

\input{texparts/python/coverage_income_scatter}



